\section{Huggett (1993)}\label{sec:huggett}

In Huggett's exchange economy households trade a bond in zero net supply at price $q$; the rate is
$r=1/q-1$, the borrowing limit $\underline a<0$ is a credit limit, and the equilibrium condition is
$K(\mu)=0$. His calibration is $\beta=0.99322$, $\sigma\in\{1.5,3\}$, a two-state endowment, and
credit limits $\underline a\in\{-2,-4,-6,-8\}$ in units of the high endowment. Two things change.

\paragraph{Single crossing needs strictness.}
Against a constant supply, a monotone $K(r)$ gives nothing: two rates with $K=0$ are compatible
with $K\equiv0$ between them. What is needed is that $K$ be strictly increasing, and this follows
if the higher-rate policy dominates the lower-rate one everywhere and strictly on a set the
lower-rate stationary distribution charges (\lean{aggregateCapital\_lt\_of\_policy\_lt},
\lean{huggett\_equilibriumRate\_unique}). The strict form of Theorem 1 holds for creditors: at
$a>0$ with an interior choice, and $r_1<r_2$, $g_1(a,z)<g_2(a,z)$ under \eqref{eq:EC}
(\lean{policy\_lt\_policy\_withRate\_of\_marginal\_at}).

\paragraph{Debtors lose cash on hand when the rate rises.}
For $a<0$, the step $c_2(a)>c_1(a)+(r_2-r_1)a$ has a negative last term, and the chain gives only
\begin{equation}\label{eq:loss}
  g_1(a,z)\ \le\ g_2(a,z)+(r_2-r_1)\max\{0,-a\}:
\end{equation}
saving can fall with the rate for a debtor, but by no more than the extra interest on the debt
(\lean{policy\_le\_policy\_add\_withRate\_of\_marginal\_at}; the slack version is
\lean{policy\_le\_policy\_add\_withRate\_of\_slack\_at}). For $a\ge0$ the loss vanishes and
Theorem 1 is as before.

\paragraph{Carrying the loss through the coupling.}
The coupling of Section~\ref{sec:coupling} needs exact dominance. A version that tolerates a shift
replaces first-order dominance by \emph{dominance up to $\delta$}: $\int h\,d\mu\le\int h\,d\nu+\delta L$
for every test function increasing in assets and $L$-Lipschitz in assets. One period of the
operator turns a shift $\delta$ into $\delta K+\gamma$, where $\gamma$ is the pointwise policy gap
\eqref{eq:loss} and $K$ is a Lipschitz constant of the higher-rate policy in assets; iterating and
passing to the limit,
\begin{equation}\label{eq:shift}
  K(\mu_1)\ \le\ K(\mu_2)+\frac{(r_2-r_1)\,|\underline a|}{1-K}\qquad\text{if }K<1
\end{equation}
(\lean{DominatesShift}, \lean{dominatesShift\_iterate},
\lean{huggett\_aggregateCapital\_le\_add}). Here $K$ is the marginal propensity to save out of
assets. For CRRA the policy's slope in assets is $R(1-\mathrm{MPC})$, and asymptotically the MPC is
$\kappa$, giving $K=R(1-\kappa)=\Thorn=(\beta R)^{1/\sigma}<1$: impatience, with almost no room. At
Huggett's numbers $1-\Thorn_2\approx0.00067$, so the amplification in \eqref{eq:shift} is about
$1500$.

\paragraph{The propensity-to-save bound cannot be had on a bounded region.}
It is tempting to derive the MPC floor from Carroll--Kimball concavity of the consumption function
together with the minimal-MPC level bound $c(x)\ge\kappa x$. On a ray this works: a concave function
above the line $\kappa x$ has every secant slope at least $\kappa$, by one chord far to the right
(\lean{secant\_ge\_of\_concaveOn\_Ici}, with the chord chosen explicitly). On the region
$[\underline a,\bar a]$ it does not: the constant function at the level of the bound at the cap is
concave, lies above the bound, and is flat (\lean{exists\_concaveOn\_ge\_flat}). The MPC floor is an
asymptotic property of the uncapped problem, and the chord at the cap says nothing because there
the level bound is far below actual consumption. So \eqref{eq:shift} is restated with the floor as a
hypothesis in the form it is used: a consumption slope floor $(R_2-K)(b-a)\le c_2(b)-c_2(a)$ is
exactly the Lipschitz bound $K$ for the policy (\lean{policy\_lipschitz\_of\_consumption\_slope}),
and the CRRA form takes $K=\Thorn_2$ and the floor $\kappa_2R_2(b-a)\le c_2(b)-c_2(a)$
(\lean{crra\_huggett\_aggregateCapital\_le\_add\_of\_slack}).

\paragraph{Where Huggett stands.}
Uniqueness reduces to a quantitative race: the creditors' strict gain in \lean{huggett\_equilibriumRate\_unique}
against the debtors' loss bounded by \eqref{eq:shift}. We have the sign of the gain and no lower
bound on its size; every attempt at one runs into a pivot condition of the type \eqref{eq:EC}. And
the MPC floor on the region is the same asymptotic structure that blocks the slack condition at
$\mu>1$. Both are open.
